% ECONOMICS_PROBLEM: {"id": "I25-601", "title": "Solving the learning crisis at scale", "jel_code": "I25", "source_id": "OP-0036"}
% !TeX program = xelatex
\documentclass[11pt,a4paper]{article}
\usepackage[margin=23mm]{geometry}
\usepackage{fontspec}
\setmainfont{Latin Modern Roman}
\usepackage{amsmath,amssymb,mathtools}
\usepackage{xcolor,enumitem,booktabs,longtable,array,xurl,hyperref}
\definecolor{muted}{HTML}{53636C}
\hypersetup{colorlinks=true,urlcolor=blue,linkcolor=blue}
\setlength{\parindent}{0pt}
\setlength{\parskip}{5pt}
\newcommand{\R}{\mathbb R}
\newcommand{\E}{\mathbb E}
\newcommand{\N}{\mathbb N}
\newcommand{\ind}{\mathbf 1}
\newcommand{\Var}{\operatorname{Var}}
\newcommand{\Cov}{\operatorname{Cov}}
\newcommand{\supp}{\operatorname{supp}}
\DeclareMathOperator*{\argmin}{arg\,min}
\DeclareMathOperator*{\argmax}{arg\,max}
\newcommand{\entrymeta}[3]{{\sffamily\small\color{muted}\textbf{\texttt{#1}}\quad|\quad #2\par #3\par}}
\newcommand{\Problemref}[1]{\texttt{#1}}
\newcommand{\JELref}[2]{\texttt{#2} (JEL \texttt{#1})}
\begin{document}
\section*{Solving the learning crisis at scale}
\textbf{Problem ID:} \texttt{I25-601}\quad \textbf{JEL:} \texttt{I25}\par
% BEGIN ECONOMICS STATEMENT
\label{problem:OP-0036}
\label{audited:ECON-044}
\label{audited:ECON-045}

\label{atlas:prob:D6}
\noindent\textbf{Original atlas ID:} D6 (entry 36). \textbf{Classification:} Editorial implementation, coverage, and learning-policy problem.

\noindent\textbf{Original atlas question:} Which pedagogy, incentives and technologies produce durable learning when implemented by normal school systems?

\subsection*{Definitions and assumptions}
Fix a school system, baseline pupils, a horizon $H$, and a common outcome scale. Policy $\pi=(p,a,s)$ specifies pedagogy or technology $p$, implementation architecture $a$ (staffing, training, supervision, and incentives), and coverage $s$. Treatment versions include time on task and language. Let $Y_{it}(\pi)$ be learning measured by a prespecified assessment with invariant interpretation across treatment groups, and let $C(\pi)$ be discounted real system cost. Follow baseline pupils across schools and nonattendance. Permit peer, teacher-labor-market, and administrative spillovers; state an exposure model or a system-level equilibrium. Assignment must occur at a level that identifies the intended policy comparison.
\subsection*{Formal research question}
For feasible pairs $\pi,\pi_0$, identify or bound
\[
 \tau_t(\pi,\pi_0)=\mathbb E[Y_{it}(\pi)-Y_{it}(\pi_0)],
 \qquad
 N_{\lambda}(\pi)=\tau_H(\pi,\pi_0)
                         -\lambda[C(\pi)-C(\pi_0)],
\]
where $\lambda\geq0$ is an explicitly chosen conversion between learning units and cost per baseline pupil, using the same population normalization for both. Determine which architectures achieve positive, persistent learning effects under ordinary system staffing and budgets. In factorial designs with supported treatment combinations, identify pedagogy-architecture interactions and how they change with coverage; otherwise bound unsupported national-coverage effects using stated model restrictions.
\subsection*{Interpretation and scope}
A successful classroom intervention with exceptional personnel does not define the same treatment as a system rollout. Architecture is part of the policy, not nuisance noise. Incentives can change teacher effort, assignment, test preparation, and reporting, so independent measurement matters. Treatment-dependent attendance or testing creates selection; follow-up or missing-outcome bounds are required. Standardized test scores can lose comparability if tests or reference distributions change, and short-term test improvement need not imply durable learning. Cost-effectiveness comparisons must use common outcome and cost units and include implementation demands. The criterion above is a chosen policy objective, not an assumption-free welfare ranking. General-equilibrium effects on teacher recruitment, school choice, and household inputs require separate identification or structural modeling.
\subsection*{Source audit and status}
[S1] verifies Indian education experiments and [S2] a 2020 school-management working paper. The intended identity of ``Mbiti et al. (2019)'' is unresolved: both the six-author input-incentive paper [S3] and the teacher-pay working paper [S4] are verified candidates. Added [S5] is a 2025 scale-adaptation experiment. Thus scaling is already directly studied; the editorial question concerns supported architecture and coverage contrasts, not a universally unexamined intervention class.

\subsection*{References}
\begin{enumerate}[label=\textnormal{[S\arabic*]},leftmargin=*]
\item Abhijit V. Banerjee, Shawn Cole, Esther Duflo, and Leigh Linden (2007). \emph{Remedying Education: Evidence from Two Randomized Experiments in India}. The Quarterly Journal of Economics 122(3), 1235--1264. \url{https://doi.org/10.1162/qjec.122.3.1235}. \textit{Locator:} Primary publisher, working-paper, or author record: bibliographic information and abstract; full-text theorem verification is not claimed. \textit{Role:} Randomized remedial and computer-assisted education evidence.
\item Karthik Muralidharan and Abhijeet Singh (2020). \emph{Improving Public Sector Management at Scale? Experimental Evidence on School Governance India}. NBER Working Paper 28129. \url{https://www.nber.org/papers/w28129}. \textit{Locator:} Primary publisher, working-paper, or author record: bibliographic information and abstract; full-text theorem verification is not claimed. \textit{Role:} Large-scale governance evaluation; verified as a working paper.
\item Isaac Mbiti, Karthik Muralidharan, Mauricio Romero, Youdi Schipper, Constantine Manda, and Rakesh Rajani (2019). \emph{Inputs, Incentives, and Complementarities in Education: Experimental Evidence from Tanzania}. The Quarterly Journal of Economics 134(3), 1627--1673. \url{https://academic.oup.com/qje/article/134/3/1627/5479257}. \textit{Locator:} Primary publisher, working-paper, or author record: bibliographic information and abstract; full-text theorem verification is not claimed. \textit{Role:} Verified thematic candidate on input-incentive complementarities. Abbreviated atlas citation has no title or link and does not uniquely distinguish this paper from the second 2019 Mbiti-led paper supplied below; existence is verified but intended identity remains unresolved.
\item Isaac Mbiti, Mauricio Romero, and Youdi Schipper (2019). \emph{Designing Effective Teacher Performance Pay Programs: Experimental Evidence from Tanzania}. NBER Working Paper 25903. \url{https://www.nber.org/papers/w25903}. \textit{Locator:} Primary publisher, working-paper, or author record: bibliographic information and abstract; full-text theorem verification is not claimed. \textit{Role:} Second verified candidate for the ambiguous shorthand; retained as a working-paper version.
\item Karthik Muralidharan and Abhijeet Singh (2025). \emph{Adapting for Scale: Experimental Evidence on Technology-Aided Instruction in India}. NBER Working Paper 34205. \url{https://www.nber.org/papers/w34205}. \textit{Locator:} Primary publisher, working-paper, or author record: bibliographic information and abstract; full-text theorem verification is not claimed. \textit{Role:} More recent experiment explicitly addresses scalability of instruction.
\end{enumerate}

\subsection*{Integrated refinement: ECON-044}
{\sffamily\small\color{muted}Scaling computer-assisted learning with existing school staff\par}
This refinement adopts the additional source conventions
(page~\pageref{audited:conventions}), subject to its local restrictions.
\textbf{Ordinary staff versus externally supplied implementers.}
For coverage $s$ and staffing $e\in\{0,1\}$, where $e=1$ is regular school staff,
define $Y_i(s,e)$ and full delivery cost $C(s,e)$. Identify
\[
\tau^{\rm staff}=\mathbb E[Y_i(s_1,1)-Y_i(s_0,1)],\qquad
\Gamma=\tau^{\rm staff}
       -\mathbb E[Y_i(s_1,0)-Y_i(s_0,0)].
\]
Use the same baseline pupils, assessment and horizon. Coverage may change
supervision, teacher time and quality in treated and untreated schools; specify
these exposures and assignment support. Include training, monitoring, staff
opportunity costs and maintenance. A pilot using exceptional personnel need not
identify the regular system's scale-up effect or implementation interaction.
\subsection*{Supplied source audit for ECON-044}
Local [S1], [S2], etc. in this audit and the references immediately below
belong to ECON-044, independently of the original entry's reference numbering.
[S1]'s evidence-gap chapter directly discusses scaling computer-assisted learning through existing teachers; [S2] corroborates the topic. The saturation/staff contrast is an editorial specification.
\subsection*{References for integrated refinement ECON-044}
{\footnotesize\raggedright
\par\noindent\textbf{[S1]} Abhijeet Singh, Laia Navarro-Sola, and Philip Oreopoulos (2025). \emph{Education Technology}. VoxDevLit 20(1), December 2025.\par \textit{Verified locator / source role:} Conclusion/evidence-gap chapter at the linked primary review URL, inspected for the agenda. See the audit conclusion for distinctions between direct agenda support and editorial extensions.\par \url{https://voxdev.org/voxdevlit/education-technology/edtech-evidence-gaps}\par\smallskip
\par\noindent\textbf{[S2]} Oliver Hanney / VoxDev (living compilation). \emph{Evidence gaps: Key unanswered questions in development economics}. Accessed 8 October 2026. The page currently covers 20 topics; the ``16-topics URL is a historical slug.\par \textit{Verified locator / source role:} Topic heading: Education Technology. Supports the broad research agenda; this is not a verbatim quotation or an attribution of the displayed mathematical target.\par \url{https://voxdev.org/topic/evidence-gaps-key-unanswered-questions-across-16-topics-development-economics}\par\smallskip
}

\subsection*{Integrated refinement: ECON-045}
{\sffamily\small\color{muted}Separating instructional practices from bundled education technology\par}
This refinement adopts the additional source conventions
(page~\pageref{audited:conventions}), subject to its local restrictions.
\textbf{Separating technology from instructional practice.}
For supported factorial assignments of technology $T$ and practice $P$,
identify $\tau_T(p)=\mathbb E[Y_i(1,p)-Y_i(0,p)]$,
$\tau_P(t)=\mathbb E[Y_i(t,1)-Y_i(t,0)]$, and
$\iota=\tau_T(1)-\tau_T(0)$. Specify hardware, software, feedback practice,
teacher preparation, baseline population and common delayed assessments.
Use consistent assignment or a justified exposure design for spillovers and
account for missing pupils. Bundle effects alone do not identify components;
label unsupported combinations as unidentified rather than infer them from a
single technology rollout. Assess persistence and comparable full resource
costs rather than immediate assisted scores alone.
\subsection*{Supplied source audit for ECON-045}
Local [S1], [S2], etc. in this audit and the references immediately below
belong to ECON-045, independently of the original entry's reference numbering.
[S1] discusses promising instructional technologies and delivery practices, with [S2] as corroboration. This technology--pedagogy factorial decomposition is an editorial extension, not a named open conjecture in the review.
\subsection*{References for integrated refinement ECON-045}
{\footnotesize\raggedright
\par\noindent\textbf{[S1]} Abhijeet Singh, Laia Navarro-Sola, and Philip Oreopoulos (2025). \emph{Education Technology}. VoxDevLit 20(1), December 2025.\par \textit{Verified locator / source role:} Conclusion/evidence-gap chapter at the linked primary review URL, inspected for the agenda. See the audit conclusion for distinctions between direct agenda support and editorial extensions.\par \url{https://voxdev.org/voxdevlit/education-technology/edtech-evidence-gaps}\par\smallskip
\par\noindent\textbf{[S2]} Oliver Hanney / VoxDev (living compilation). \emph{Evidence gaps: Key unanswered questions in development economics}. Accessed 8 October 2026. The page currently covers 20 topics; the ``16-topics URL is a historical slug.\par \textit{Verified locator / source role:} Topic heading: Education Technology. Supports the broad research agenda; this is not a verbatim quotation or an attribution of the displayed mathematical target.\par \url{https://voxdev.org/topic/evidence-gaps-key-unanswered-questions-across-16-topics-development-economics}\par\smallskip
}

\subsection*{Common conventions: Source mathematical conventions}

Unless an entry states otherwise, agent and firm sets are finite. State and action spaces are measurable, strategies and outcome maps are measurable, probability laws are specified, and expectations exist for the targets under discussion. Infinite-horizon discount factors lie in $(0,1)$ and discounted objectives are finite. Norms, tolerances and complexity input sizes must be specified for computational comparisons. Constants or primitives that appear locally are inputs, not empirical facts asserted by this catalogue.

\textbf{Identification.} A statistical model $m\in\mathcal M$ implies an observable law $P_m$ and target $\theta(m)$. At law $P$, its identified set is
\[
\Theta(P;\mathcal M)=\{\theta(m):m\in\mathcal M,\ P_m=P\}.
\]
Point identification holds when this set is a singleton, or equivalently when $P_{m_1}=P_{m_2}$ implies $\theta(m_1)=\theta(m_2)$ for all compatible models. Sharp bounds mean bounds attained, or approached when the boundary is not attained, by compatible models. Writing this set is a definition, not a solution: the substantive task is to establish informative restrictions and derive or estimate the set. An unrestricted-confounding model may give only trivial bounds.

\textbf{Inference.} Identification, estimability and valid uncertainty quantification are separate questions. Fix $\alpha\in(0,1)$. A confidence procedure $C_n$ over a declared law class $\mathcal P$ has asymptotic uniform coverage at level $1-\alpha$ when
\[
\liminf_{n\to\infty}\inf_{P\in\mathcal P}\Pr_P\{\theta(P)\in C_n\}\geq1-\alpha.
\]
For set-identified targets the coverage event must be defined separately, for example coverage of the full identified set. If an entry uses identification-indexed models instead of uniquely indexed laws, the same coverage convention is applied over those models. Pointwise limits do not establish uniform validity, and asymptotic validity does not guarantee finite-sample precision. Sample sizes, dependence structures and weak-identification sequences are part of each application.

\textbf{Counterfactuals.} $Y_i(a)$ is a potential outcome under a specified intervention, including its timing, implementation and versions. It is not $Y_i$ conditional on voluntarily choosing $a$. A full intervention vector permits interference; shorthand scalar treatment notation does not silently impose no interference. Assignment, selection, attrition, anticipation and measurement must be specified. A claim about an instrument's local effect is not automatically a population policy effect.

\textbf{Economic equilibrium.} A counterfactual model includes best responses, constraints, feasibility, market clearing, state transitions and expectations consistent with the stated information. When several equilibria exist, either a declared selector is an input or the target is set-valued. Comparisons across policy regimes must use the stated selection convention. Reduced-form relationships are not presumed invariant after an equilibrium-changing intervention.

\textbf{Optimization and welfare.} Optimization is over the stated feasible set. If attainment is not established, $\sup$ or $\inf$ and $\varepsilon$-optimal choices replace an asserted maximizer or minimizer. For example, with value $V=\sup_{a\in\mathcal A}W(a)<\infty$, an $\varepsilon$-optimal set is $\{a\in\mathcal A:W(a)\geq V-\varepsilon\}$ for $\varepsilon>0$. Benefits, costs, risk and health or environmental outcomes can be aggregated only using declared units and conversion weights. Interpersonal utility weights, the treatment of fiscal transfers and foreign profits, discounting, and the behavioral welfare criterion are explicit inputs. A comparative-static derivative additionally requires existence and appropriate differentiability of the selected counterfactual.

\textbf{Sources.} Local labels [S1], [S2], and so forth restart in every question. Locators identify the relevant abstract, model, section or result at the level actually checked; a bibliographic or abstract-level verification is not represented as inspection of an entire proof. Working papers are distinguished from later publications and changing versions. Source summaries are paraphrased. All URL checks and editorial formulations refer to the audit date stated above.

\subsection*{Common conventions: Additional-source status and mathematical conventions}

\label{audited:conventions}

Of the 100 supplied ECON entries, 65 distinct problems appear here; 32 close
matches refine earlier entries, and three duplicate questions map to existing
entries. Appendix~\ref{app:audited-crosswalk} lists every destination.
The attachment's P/R/S/U distinctions and source audits are retained. Its
precise mathematical formalizations are editorial unless the individual audit
says otherwise. Candidate subsidy targets ECON-004--006 retain their U flags;
the resolved approval-core question ECON-007 updates OP-0202 above.
The following supplied conventions govern both this part and the attributed
ECON refinements elsewhere in the catalogue, subject to local restrictions.

\textbf{Parameterized questions.} A problem family lists inputs that must be fixed before an instance is evaluated: population, horizon, policies, information, data law, model class and welfare weights. These are required choices, not quantities the citations are assumed to specify. Undefined applications should not be treated as identified estimands.

\textbf{Indices and integrability.} Sets of agents, firms and regions are finite unless a population measure is explicitly used. \(i,j\) index units, \(r\) regions, \(t\) dates and \(h\) horizons. \(\mathbb E\) and \(\Pr\) refer to the declared population and law; \(\mathbf1\{A\}\) is an indicator. Every displayed expectation and discounted sum needs existence in the stated domain. Infinite horizons use a discount in \((0,1)\) and a convergence condition; finite horizons may use discount one.

\textbf{Optimization and existence.} A displayed \(\max\), \(\min\), or \(\arg\max\) is conditional on attainment. For finite feasible menus this is immediate; general spaces need continuity and compactness/coercivity or another existence argument. Without attainment, the target is the supremum/infimum and arbitrarily near-optimal feasible choices. \(\inf\varnothing=+\infty\). Empty feasibility and an unidentified target are different issues.

\textbf{Potential outcomes.} \(Y_i(z)\) is the outcome under a specified intervention, not the observed conditional mean among people choosing \(z\). In binary comparisons \(\Delta Y_i=Y_i(1)-Y_i(0)\). Specify assignment, treatment versions, compliance, information and observation horizon. Interference is allowed under a full policy vector; compressed exposure notation needs a justified mapping. No interference, consistency and overlap are substantive assumptions, not automatic consequences of notation.

\textbf{Populations and observation.} Fix baseline eligibility and population weights. Conditioning on post-policy employment, survival, relocation or program uptake can change the population being compared. Death is not ordinary loss to follow-up; outcomes truncated by death need a composite convention, joint survival target, or explicitly defined survivor stratum. Fertility-changing interventions need a population functional rather than an unexplained fixed descendant cohort. Measurement and missingness models must be supplied.

\textbf{Identification.} A structural model \(m\in\mathcal M\) implies observable law \(P_m\) and target \(\theta(m)\). Identification means
\[
 P_{m_1}=P_{m_2}\ \Longrightarrow\ \theta(m_1)=\theta(m_2)
 \quad\text{for all }m_1,m_2\in\mathcal M .
\]
Without identification, the target is
\[
 \Theta(P;\mathcal M)=\{\theta(m):m\in\mathcal M,\ P_m=P\}.
\]
Writing \(\theta(P)\) before establishing identification can hide incompatible latent models. Confidence sets additionally need a sampling law, confidence level, asymptotic sequence and uniformity class. Point identification, coverage and informative precision are separate properties.

\textbf{Mediation.} Total effects of randomized assignment are distinct from cross-world natural indirect effects. Belief/effort-channel effects require a meaningful mediator intervention and additional measurement, exclusion or mediation assumptions. Randomization of the initial policy alone does not supply those assumptions. Report bounds or interventional alternatives when the stronger target is unsupported.

\textbf{Equilibrium.} A counterfactual \(W(z)\), \(p(z)\), or \(\mu_t^z\) requires individual optimization, feasibility, government/resource budgets, market clearing and state-transition laws. State equilibrium existence and selection, or report ranges over compatible equilibria. Initial-state integration includes subsequent endogenous transitions; current-state integration must use the policy-dependent distribution. Reduced-form invariance after policy is not assumed.

\textbf{Welfare accounts.} Scores, utility, health, dollars and ecological quantities require explicit conversion before addition. Fees, taxes, subsidies, wages and extortion payments are transfers between parties; distributional weights can make them welfare relevant, but their face value is not automatically a resource loss. State ownership, geographic boundaries, foreign-income weights and fiscal finance. Do not add profits to household welfare already including dividends, or count land capitalization and the underlying benefit twice. A benefit-cost expression must distinguish gross benefits from net welfare and subtract every real cost exactly once.

\textbf{Derivatives and risk.} Log derivatives require positive arguments; local responses require admissible perturbations and specified equilibrium branches. Differentiation under expectations needs regularity/domination, and boundaries may allow only one-sided derivatives. Risk uses a specified law; ambiguity uses a declared nonempty law set on a common history space. Adaptive policies must be nonanticipative. A robust criterion is an editorial choice unless attributed explicitly.

\textbf{Scope overlaps.} Allocation, collective choice, contracts, household bargaining and coordinated policy overlap economics with optimization and game theory. They are retained because they appeared in the original catalogue; no claim of a strictly game-theory-free selection is made.
% END ECONOMICS STATEMENT
\end{document}
